% Borcherds singular theta lift for the Igusa square-root normalization.
% Standalone Vol III note.

\documentclass[11pt]{amsart}
\usepackage[localtheorems]{raeez-math-template}

\newtheorem{theorem}{Theorem}[section]
\newtheorem{proposition}[theorem]{Proposition}
\newtheorem{lemma}[theorem]{Lemma}
\newtheorem{corollary}[theorem]{Corollary}
\newtheorem{definition}[theorem]{Definition}
\newtheorem{remark}[theorem]{Remark}
\newtheorem{obligation}[theorem]{Proof obligation}

\newcommand{\kch}{\kappa_{\mathrm{ch}}}
\newcommand{\kBKM}{\kappa_{\mathrm{BKM}}}
\newcommand{\ZZ}{\mathbb{Z}}
\newcommand{\QQ}{\mathbb{Q}}
\newcommand{\RR}{\mathbb{R}}
\newcommand{\CC}{\mathbb{C}}
\newcommand{\HH}{\mathbb{H}}
\newcommand{\Sp}{\mathrm{Sp}}
\newcommand{\OO}{\mathrm{O}}
\newcommand{\Oplus}{\mathrm{O}^{+}}
\newcommand{\Mp}{\mathrm{Mp}}
\newcommand{\Sieg}{\mathbb{H}_2}
\newcommand{\fg}{\mathfrak{g}}
\newcommand{\gDelta}{\mathfrak{g}_{\Delta_5}}
\newcommand{\phij}{\phi_{0,1}}
\newcommand{\Delfive}{\Delta_5}
\newcommand{\PhiTen}{\Phi_{10}}
\newcommand{\Weil}{\rho}
\newcommand{\Bor}{\operatorname{Borch}_{\theta}}
\newcommand{\wt}{\operatorname{wt}}
\newcommand{\divisor}{\operatorname{div}}

\begin{document}

\title{The singular theta lift of \texorpdfstring{$\phij$}{phi01}
and the Borcherds product \texorpdfstring{$\Delfive$}{Delta5}}
\author{}
\date{}
\maketitle

\begin{abstract}
The weight-five Igusa square root \(\Delfive\) is the Borcherds
automorphic product attached to the weight-zero index-one weak Jacobi
form \(\phij\), after fixing the lattice \(2U\oplus\langle -2\rangle\),
the Eichler--Zagier theta decomposition, the Borcherds product
normalization, and the theta-constant normalization of \(\Delfive\).
In this normalization the monic product is
\[
 D_5(Z)=64^{-1}\Delfive(Z),
 \qquad
 \operatorname{den}(\gDelta)=D_5(2Z)=64^{-1}\Delfive(2Z).
\]
The Borcherds weight theorem then gives
\[
 \kBKM(\Delfive; \phij,2U\oplus\langle -2\rangle,D_5)
 =\frac{c(0,0)}2=\frac{10}{2}=5.
\]
The formula is a theorem only after the input Jacobi form, lattice,
group or cover, character, and product normalization have been named.
\end{abstract}

\tableofcontents

\section{Lattice, Jacobi input, and Weil representation}

\begin{definition}[The orthogonal model]
Let
\[
 M=2U\oplus\langle -2\rangle
\]
with \(U\) the hyperbolic plane.  Then \(M\) has signature \((2,3)\)
and discriminant group \(A_M=M^\vee/M\simeq \ZZ/2\ZZ\).  The stable
orthogonal group of \(M\) is the orthogonal model for the genus-two
paramodular group at index \(1\):
\[
 \widetilde{\OO}^{+}(M)\simeq \Sp_4(\ZZ)/\{\pm I\}.
\]
The sign-reversed lattice
\[
 L=-M\simeq U\oplus U\oplus\langle 2\rangle
\]
has signature \((3,2)\).  Statements below use \(M\), because
Borcherds' theorem is stated for signature \((2,b^-)\).  The \(L\)
formulation is obtained by changing the sign convention on the
quadratic form and the associated Weil representation.
\end{definition}

\begin{definition}[The Jacobi input]
Let
\[
 \phij(\tau,z)=\sum_{n\geq 0,\ell\in\ZZ} f(n,\ell)q^n r^\ell
 \in J^{\mathrm{wk}}_{0,1},
 \qquad q=e^{2\pi i\tau},\quad r=e^{2\pi i z},
\]
be the standard weight-zero index-one weak Jacobi form
\[
 \phij(\tau,z)
 =r^{-1}+10+r
 +q(10r^{-2}-64r^{-1}+108-64r+10r^2)+O(q^2).
\]
Thus
\[
 f(0,-1)=1,\qquad f(0,0)=10,\qquad f(0,1)=1.
\]
\end{definition}

\begin{proposition}[Eichler--Zagier conversion]
\label{prop:ez-conversion}
The theta decomposition
\[
 \phij(\tau,z)
 =h_0(\tau)\vartheta_{0,1}(\tau,z)
  +h_1(\tau)\vartheta_{1,1}(\tau,z)
\]
defines a vector-valued weakly holomorphic modular form
\[
 F_{\phij}=h_0\,\mathfrak e_0+h_1\,\mathfrak e_1
 \in M^!_{-1/2}(\Weil_M)
\]
for the Weil representation of \(A_M\simeq\ZZ/2\ZZ\), after the
standard identification of the index-one Jacobi discriminant form
with the discriminant form of \(\langle -2\rangle\).  Its principal
terms are
\[
 F_{\phij}(\tau)
 =q^{-1/4}\mathfrak e_1+10\,\mathfrak e_0+O(q^{3/4}).
\]
\end{proposition}

\begin{proof}
Eichler--Zagier identify weak Jacobi forms of weight \(k\) and index
\(m\) with vector-valued weakly holomorphic modular forms of weight
\(k-1/2\) for the Weil representation of the rank-one lattice
\(\langle 2m\rangle\), with dual discriminant convention if the
orthogonal lattice uses \(\langle -2m\rangle\).  For \(k=0\) and
\(m=1\), this gives weight \(-1/2\).  The \(q^0r^0\)-coefficient of
\(\phij\) is \(10\), and the \(q^0r^{\pm1}\)-terms give the single
polar vector-valued term \(q^{-1/4}\mathfrak e_1\).
\end{proof}

\begin{remark}[No weight mismatch]
\label{rem:no-weight-mismatch}
Borcherds' product theorem for a lattice of signature \((2,b^-)\)
uses vector-valued input of weight \(1-b^-/2\).  Here \(b^-=3\), so
the required weight is \(1-3/2=-1/2\).  The Eichler--Zagier image
of \(\phij\) has exactly this weight.  No auxiliary induction from
weight \(-1/2\) to weight \(-3/2\) is present in the \(\Delfive\)
construction.
\end{remark}

\section{The Borcherds theorem in the required normalization}

\begin{theorem}[Borcherds singular theta product]
\label{thm:borcherds-product}
Let \(M\) be an even lattice of signature \((2,b^-)\), and let
\[
 F(\tau)=\sum_{\gamma\in A_M}\sum_{m\in\QQ}
 c_\gamma(m)q^m\mathfrak e_\gamma
 \in M^!_{1-b^-/2}(\Weil_M)
\]
be holomorphic on \(\HH\), meromorphic at the cusp, and integral in
all non-positive Fourier coefficients.  Borcherds' regularized theta
integral, defined as the constant term at \(s=0\) of the analytically
continued truncated integral
\[
 \int_{\mathcal F}
 \bigl\langle F(\tau),\overline{\Theta_M(\tau,Z)}\bigr\rangle
 y^{-s}\,\frac{dx\,dy}{y^2},
\]
produces a meromorphic automorphic form
\[
 \Psi_M(Z,F)
\]
on the type IV domain for \(M\).  Its weight is
\[
 \wt(\Psi_M)=\frac{c_0(0)}{2}.
\]
Its divisor is supported on rational quadratic divisors
\(\lambda^\perp\) with \(\lambda^2<0\), with order
\[
 \sum_{\substack{x>0\\x\lambda\in M^\vee}}
 c_{x\lambda}(x^2\lambda^2/2).
\]
At a primitive zero-dimensional cusp and in a Weyl chamber \(W\), it
has product expansion
\[
 \Psi_M(Z,F)
 =C\,e^{2\pi i(\rho,Z)}
 \prod_{\substack{\lambda\in K^\vee\\(\lambda,W)>0}}
 \prod_{\substack{\delta\in A_M\\\delta|K=\lambda}}
 \bigl(1-e^{2\pi i((\lambda,Z)+(\delta,z'))}\bigr)^{
 c_\delta(\lambda^2/2)},
\]
where \(K=(\ZZ z+\ZZ z')^\perp/(\ZZ z+\ZZ z')\), \(\rho\) is the
Weyl vector, and \(C\) is the cusp constant fixed by the chosen
normalization.
\end{theorem}

\begin{proof}
This is Borcherds, \emph{Automorphic forms with singularities on
Grassmannians}, Theorem 13.3, with the regularized theta integral
defined in Sections 6 and 7 of that paper.  The theorem states the
weight, divisor order, logarithmic relation with the regularized
theta integral, and product expansion at each primitive norm-zero
cusp.
\end{proof}

\begin{corollary}[The theorem applies to \(\phij\)]
\label{cor:borch-applies-phi01}
For \(M=2U\oplus\langle -2\rangle\), the form \(F_{\phij}\) of
Proposition~\ref{prop:ez-conversion} satisfies all hypotheses of
Theorem~\ref{thm:borcherds-product}.  Therefore
\[
 \wt\bigl(\Psi_M(Z,F_{\phij})\bigr)=\frac{10}{2}=5,
\]
and the only primitive polar divisor is the Humbert divisor attached
to the class \(\mathfrak e_1 q^{-1/4}\).
\end{corollary}

\begin{proof}
The lattice has \(b^-=3\), so the input weight is \(-1/2\), matching
Proposition~\ref{prop:ez-conversion}.  The non-positive coefficients
visible at the cusp are integral:
\[
 c_1(-1/4)=1,\qquad c_0(0)=10.
\]
Borcherds' divisor formula assigns the polar term \(q^{-1/4}
\mathfrak e_1\) to the corresponding rational quadratic divisor, and
the constant term \(c_0(0)=10\) to the automorphic weight.
\end{proof}

\section{The Igusa form as the output of the lift}

\begin{theorem}[Gritsenko--Nikulin product for \(\Delfive\)]
\label{thm:delta5-output}
In the theta-constant normalization used here, the Borcherds product
of \(F_{\phij}\) is the monic Igusa square root
\[
 D_5(Z)=\Psi_M(Z,F_{\phij})=64^{-1}\Delfive(Z).
\]
Equivalently, in the even denominator-lattice variables used for
\(\gDelta\),
\[
 \operatorname{den}(\gDelta)=D_5(2Z)=64^{-1}\Delfive(2Z).
\]
The product expansion is
\[
 D_5(Z)
 =q^{1/2}r^{1/2}s^{1/2}
 \prod_{(n,\ell,m)>0}
 \bigl(1-q^n r^\ell s^m\bigr)^{f(nm,\ell)}
\]
up to the fixed choice of the positive Weyl chamber, with
\(f(n,\ell)\) the Fourier coefficients of \(\phij\).
\end{theorem}

\begin{proof}
Theorem~\ref{thm:borcherds-product} gives a weight-five automorphic
product on the orthogonal model for \(\Sp_4(\ZZ)\), with divisor
the diagonal Humbert divisor of multiplicity one.  Gritsenko--Nikulin
identify the resulting product with the weight-five Igusa square root,
and Gritsenko--Clery record the same row as the first diagonal-divisor
form in their classification.  The theta-constant convention fixes
the leading coefficient
\[
 [q^{1/2}r^{1/2}s^{1/2}]\Delfive=64.
\]
The Borcherds product is monic at the cusp, so the output is
\(D_5=64^{-1}\Delfive\).  Replacing \(Z\) by \(2Z\) is the separate
denominator-lattice normalization used for the Gritsenko--Nikulin
superalgebra \(\gDelta\); it changes the displayed argument, not the
weight.
\end{proof}

\begin{corollary}[The Borcherds characteristic]
\label{cor:kBKM-delta5}
With the input and normalization of Theorem~\ref{thm:delta5-output},
\[
 \kBKM(\Delfive;\phij,M,D_5)
 =\wt(D_5)
 =\frac{c_0(0)}{2}
 =\frac{10}{2}
 =5.
\]
\end{corollary}

\begin{proof}
The constant coefficient is the \(q^0r^0\)-coefficient of \(\phij\),
equivalently the \(\mathfrak e_0q^0\)-coefficient of \(F_{\phij}\).
Borcherds' weight theorem divides this number by \(2\).  The scalar
\(64\), the substitution \(Z\mapsto 2Z\), and the square
\(\PhiTen=\Delfive^2\) do not alter this weight.  In the square
normalization the weight is \(10\), and the relevant constant term is
twice the one used for the square-root denominator.
\end{proof}

\begin{remark}[Three nearby numbers]
\label{rem:three-nearby-numbers}
The number \(5\) in Corollary~\ref{cor:kBKM-delta5} is the Borcherds
weight.  The value \(K^{\kch}=8\), when used in the Mukai-Heisenberg
lane, is a separate Heegner or Mukai intersection invariant and is not
computed from \(c_0(0)\).  The square \(\PhiTen=\Delfive^2\) has weight
\(10\) and belongs to the scalar Igusa-square convention.  None of
these three quantities may be substituted for the others.
\end{remark}

\section{Separation from adjacent product families}

\begin{proposition}[Four non-interchangeable constructions]
\label{prop:separations}
The following objects share Borcherds-product language but are distinct.
\begin{enumerate}
\item The \(\Delfive\) product above uses \(M=2U\oplus\langle -2\rangle\),
the weak Jacobi form \(\phij\in J^{\mathrm{wk}}_{0,1}\), the Maass
character of the weight-five Igusa square root, and the monic
normalization \(D_5=64^{-1}\Delfive\).
\item The Gritsenko--Clery diagonal-divisor catalogue classifies eight
genus-two modular forms with simplest diagonal divisor.  The first row
is \(\Delfive\); the remaining rows require their own paramodular
group or cover, character, Jacobi input, divisor, and product theorem.
The catalogue is not itself a CHL theorem and is not a single universal
BKM denominator construction.
\item CHL Borcherds products are attached to specified CHL orbifold
data, twined or twisted K3 elliptic genera, and level-dependent
paramodular groups.  A CHL product may coincide with a catalogue row
only after the CHL lattice, twining class, character, Weyl chamber, and
normalization are matched.
\item The Fake Monster product is the Borcherds denominator attached
to the even unimodular Lorentzian lattice \(II_{25,1}\) and the modular
form \(\eta^{-24}\).  Its automorphic weight is \(12\).  It is not a
specialization of the \(M=2U\oplus\langle -2\rangle\) Igusa product.
\end{enumerate}
\end{proposition}

\begin{proof}
The data determining a Borcherds product are the lattice and its
discriminant form, the input modular or Jacobi form, the modular group
or cover, the character, the Weyl chamber, and the product
normalization.  The four constructions listed above differ in at least
one of these determining data.  Borcherds' weight formula can be reused
only after the corresponding data have been fixed.
\end{proof}

\section{Proof obligations for further theta-lift claims}

\begin{obligation}[General \(\kBKM=c(0)/2\) assertions]
\label{obl:general-kBKM}
A statement of the form
\[
 \kBKM(\Phi_N)=\frac{c_N(0)}2
\]
must name:
\begin{enumerate}
\item the lattice \(M_N\) and discriminant form \(A_{M_N}\);
\item the input Jacobi, vector-valued, or harmonic Maass form \(F_N\);
\item the group or metaplectic cover, including the character;
\item the cusp and Weyl chamber used for the product;
\item the product normalization defining \(\Phi_N\);
\item the theorem proving that the resulting product is reflective or
is a denominator form, if a BKM interpretation is claimed.
\end{enumerate}
Without these data the equality is only a candidate Borcherds-weight
formula, not a theorem about a named automorphic product.
\end{obligation}

\begin{obligation}[Heegner and Mukai numerical claims]
\label{obl:heegner-mukai}
Any claim identifying \(K^{\kch}=8\) with a divisor or Chern-class
calculation must specify the orthogonal Shimura variety, the Heegner
divisor, the automorphic line bundle, the Mukai or K3 moduli subvariety,
and the intersection computation.  The Borcherds theorem supplies the
divisor
\[
 \divisor(\Psi_M(Z,F))
 =\sum_{\gamma,m>0} c_\gamma(-m)Z(\gamma,m),
\]
but it does not turn \(c_0(0)=10\) into \(K^{\kch}=8\).
\end{obligation}

\begin{obligation}[Categorical lift claims]
\label{obl:categorical-lift}
The identity
\[
 \Bor(\phij)=\Delfive
\]
is an automorphic theorem.  A categorical statement relating it to a
compact \(K3\times E\) Hall or factorization-algebra construction
requires additional data: compact Hall correspondences, orientation
line and reflection-character comparison, the positive-half and
Drinfeld-double completion, a character map to the denominator product,
and a proof that the resulting bracket satisfies the Gritsenko--Nikulin
Borcherds--Kac relations.  The theta lift alone does not construct
that compact source.
\end{obligation}

\section{References used}

\begin{itemize}
\item R. E. Borcherds, \emph{Automorphic forms with singularities on
Grassmannians}, Invent. Math. 132 (1998), 491--562, Theorem 13.3.
\item V. Gritsenko and V. Nikulin, Amer. J. Math. 119 (1997), and the
companion Siegel-product paper; Example 2.4 gives the
\(\phij\) product for \(\Delfive\).
\item V. Gritsenko, \emph{Elliptic genus of Calabi--Yau manifolds and
Jacobi and Siegel modular forms}, St. Petersburg Math. J. 11 (1999),
for the paramodular lifting normalization.
\item F. Clery and V. Gritsenko, \emph{The Siegel modular forms of
genus 2 with the simplest divisor}, Proc. Lond. Math. Soc. 102
(2011), 1024--1052.
\item M. Eichler and D. Zagier, \emph{The Theory of Jacobi Forms},
Birkhauser, 1985, for the Jacobi-to-vector-valued theta decomposition.
\item The local exact coefficient check
\texttt{/Users/raeez/igusa-cusp-form/compute/verify\_square\_root.py},
which gives \(f(0,0)=10\) and the leading Igusa normalization used
above.
\end{itemize}

\end{document}
